\section{Nhóm ứng dụng di động doanh nghiệp và giáo dục đơn phân hệ}
\label{sec:app_don_phan_he}

Trên thị trường hiện nay, nhiều giải pháp ứng dụng di động phục vụ quản trị nhân sự và tác vụ hành chính đã được phát triển bởi các doanh nghiệp phần mềm hoặc được một số trường đại học ứng dụng thử nghiệm:

\begin{itemize}
    \item \textbf{Ứng dụng di động chuyên biệt cho Chấm công \& Nghỉ phép (như Base.vn, 1Office, Tanca):}
    \begin{itemize}
        \item \textit{Ưu điểm:} Giao diện di động được trau chuốt, hỗ trợ chấm công định vị GPS/Wi-Fi, tạo đơn xin nghỉ phép nhanh chóng và nhận thông báo đẩy tức thời.
        \item \textit{Hạn chế:} Đây là các giải pháp đóng gói thương mại (SaaS) được thiết kế cho khối doanh nghiệp tư nhân với quy trình hành chính đơn giản. Chúng không thể đáp ứng được các quy chuẩn quản lý nhân sự phức tạp của cơ quan nhà nước / đại học công lập (như tính toán quỹ phép năm theo thâm niên ngạch bậc, quy trình duyệt đơn 2 cấp có thẩm định số quyết định TCCB, và yêu cầu bảo mật dữ liệu lưu trữ nội bộ theo quy định của ĐHQG-HCM).
    \end{itemize}
    \item \textbf{Ứng dụng di động đơn phân hệ nội bộ tại một số trường đại học:}
    \begin{itemize}
        \item \textit{Ưu điểm:} Phục vụ nhanh một nhu cầu cấp bách (ví dụ chỉ xem thời khóa biểu, hoặc chỉ xem sổ tay sinh viên).
        \item \textit{Hạn chế:} Dẫn đến tình trạng phân mảnh ứng dụng (App Fragmentation), mỗi phân hệ là một app riêng biệt khiến cán bộ phải cài đặt nhiều ứng dụng và đăng nhập nhiều lần, trải nghiệm người dùng bị đứt gãy.
    \end{itemize}
\end{itemize}

\section{Nhóm hệ thống quản lý nhân sự và điều hành trên nền tảng Web (Web-based Portals)}
\label{sec:web_portals}

Tại hầu hết các trường đại học tại Việt Nam, bao gồm cả Trường Đại học Bách khoa -- ĐHQG-HCM, các quy trình nghiệp vụ cốt lõi đều được triển khai trên nền tảng Web Desktop:
\begin{itemize}
    \item \textbf{Hệ thống Quản lý Nhân sự (HRM Web):} Quản lý toàn bộ hồ sơ lý lịch cán bộ (11 danh mục), quá trình nâng bậc lương, khen thưởng, kỷ luật, quy trình nộp đơn nghỉ phép và đăng ký đi công tác.
    \item \textbf{Hệ thống Văn phòng điện tử (iOffice Web):} Quản lý sổ văn bản đến, văn bản đi, phân phối chỉ đạo văn bản, điều hành nhiệm vụ (Tasks) và lịch tuần công tác trường/đơn vị.
\end{itemize}

\noindent \textbf{Đánh giá ưu điểm và hạn chế của nhóm hệ thống Web:}
\begin{itemize}
    \item \textit{Ưu điểm:} Đáp ứng đầy đủ, chặt chẽ tất cả các quy chuẩn nghiệp vụ quản trị nhà trường; cơ sở dữ liệu quan hệ PostgreSQL lưu trữ tập trung, an toàn và toàn vẹn giao dịch (ACID).
    \item \textit{Hạn chế:} Giao diện Web Desktop được thiết kế với nhiều bảng biểu phức tạp, không tối ưu cho màn hình cảm ứng di động; không có cơ chế gửi thông báo đẩy (Push Notifications) về điện thoại cán bộ khi có việc khẩn; không hỗ trợ hoạt động ngoại tuyến và không tận dụng được các tính năng phần cứng của điện thoại (như điểm danh họp thời gian thực qua WebSocket).
\end{itemize}

\section{Đánh giá chung và Đề xuất giải pháp}
\label{sec:danhgiachung_dexuat}

\begin{table}[H]
    \centering
    \renewcommand{\arraystretch}{1.3}
    \begin{tabularx}{\textwidth}{|p{3.5cm}|X|X|X|}
    \hline
    \textbf{Tiêu chí Đánh giá} & \textbf{App Thương mại (SaaS)} & \textbf{Hệ thống Web Nội bộ} & \textbf{Giải pháp MyHCMUT Đề xuất} \\
    \hline
    \textbf{Phù hợp Quy chế ĐHBK} & Thấp, không tùy biến được quy trình duyệt TCCB. & Tuyệt đối (100\% đúng nghiệp vụ nhà trường). & Tuyệt đối (Tích hợp trực tiếp dữ liệu Backend trường). \\
    \hline
    \textbf{Trải nghiệm Di động (Mobile UX)} & Cao, giao diện mượt mà trên smartphone. & Kém, giao diện web bị vỡ và chữ nhỏ trên mobile. & Cao (Flutter Native UI, Material 3 Design System). \\
    \hline
    \textbf{Tính Tập trung (Unified)} & Phân mảnh, chỉ có tính năng chấm công/nghỉ phép. & Tách rời giữa các trang Web HRM và iOffice riêng. & Tập trung 100\% (HRM, iOffice, Tasks, Schedule trên 1 App). \\
    \hline
    \textbf{Bảo mật Dữ liệu Nội bộ} & Rủi ro (dữ liệu lưu trên Cloud bên thứ ba). & Cao (lưu trữ trên cụm máy chủ nội bộ trường). & Tuyệt đối (Giao tiếp bảo mật nội bộ qua Shared JWT). \\
    \hline
    \end{tabularx}
    \caption{\textit{Bảng so sánh tổng hợp các giải pháp hiện có và giải pháp đề xuất}}
    \label{tab:compare_solutions}
\end{table}

\noindent \textbf{$\rightarrow$ Đề xuất Giải pháp:}  
Từ phân tích trên, giải pháp tối ưu nhất là xây dựng **Ứng dụng Di động Cổng Giao tiếp Tập trung (Unified Mobile Portal) - MyHCMUT**. Ứng dụng này không thay thế Backend hiện hữu mà đóng vai trò là một lớp giao diện di động hiện đại (Integration Layer), kết nối an toàn với các dịch vụ Backend của Nhà trường thông qua REST API, Kafka Event Streaming và WebSocket, mang lại trải nghiệm làm việc số toàn diện cho cán bộ giảng viên.